% !TeX root = ../../金融经济学_熊和平_习题.tex

\begin{problem}
\problemno{6.1} 什么是衍生证券？衍生证券通常包括哪些种类？
\end{problem}
\begin{solution}
衍生证券是价值由标的资产、利率、汇率或指数等基础变量决定的金融合约。其支付依赖标的资产价格或状态，常见种类包括远期、期货、期权、互换，以及由这些合约组合而成的结构化产品。
\end{solution}
\begin{problem}
\problemno{6.2} 什么是积木分析法？它是如何体现无套利思想的？
\end{problem}
\begin{solution}
积木分析法把复杂衍生证券拆分为若干基本支付积木，例如零息债券、股票和状态证券，再分别定价并相加。若两种组合在每个状态支付完全相同，它们的价格必须相同；否则买入便宜组合、卖出昂贵组合即可套利。
\end{solution}
\begin{problem}
\problemno{6.3} 如何确定期权的价值？什么是内在价值？什么是时间价值？
\end{problem}
\begin{solution}
期权价值由复制或无套利上下界确定。看涨期权的内在价值为 $\max(S-K,0)$，看跌期权的内在价值为 $\max(K-S,0)$；时间价值是期权价格减去内在价值，反映到期前标的价格和波动率可能带来的额外收益。
\end{solution}
\begin{problem}
\problemno{6.4} 什么是状态指数化证券？为什么利用它能够使市场完全化？如何使市场完全化？
\end{problem}
\begin{solution}
状态指数化证券只在某一状态支付一个单位，在其他状态支付零。若有 $S$ 个状态并引入 $S$ 个相互独立的状态证券，支付矩阵就是单位矩阵，秩为 $S$，因此任意状态支付都可复制，市场完全化。实际操作是为每个不能复制的状态增加相应的状态证券。
\end{solution}
\begin{problem}
\problemno{6.5} 以二叉树模型为例，说明期权定价有哪些核心思想。
\end{problem}
\begin{solution}
核心思想包括：第一，用股票和无风险资产构造与期权相同的复制组合；第二，由一价律令期权价格等于复制成本；第三，用风险中性概率计算贴现期望支付；第四，利用向后递推从到期节点逐层求出各时点价格。风险中性概率只是定价工具，不必等于真实上涨概率。
\end{solution}
\begin{problem}
\problemno{6.6} 分别说明美式期权在支付红利和不支付红利两种情形下是否提前执行，说明理由。
\end{problem}
\begin{solution}
不支付红利的股票看涨期权通常不应提前执行：持有期权保留了上行选择权，提前执行会失去时间价值。支付红利时，除息会使股票价格下降，若红利足够大，临近除息日提前执行可能最优。美式看跌期权即使不支付红利也可能提前执行，因为立即获得执行价与股票价值之差可能超过继续持有的时间价值。
\end{solution}
\begin{problem}
\problemno{6.7} 以期权价格的上下界为例说明这些等量关系和不等量关系成立的前提条件。
\end{problem}
\begin{solution}
无套利下欧式看涨期权满足 $\max(S_0-Ke^{-rT},0)\le c\le S_0$，欧式看跌期权满足 $\max(Ke^{-rT}-S_0,0)\le p\le Ke^{-rT}$。等式需要复制或退化条件，例如到期立即执行时内在价值等于价格；不等式依赖无套利、可借贷、无交易成本和标的可交易等前提。支付红利时上下界要用标的资产净现值替换 $S_0$。
\end{solution}

\begin{problem}
\problemno{6.8} 证明期权性质 1 至性质 6。
\end{problem}
\begin{solution}
各性质均由逐状态支付比较得到。例如看涨期权支付 $(S_T-K)^+$ 随 $S_T$ 单调增加、随 $K$ 单调减少，故 $c$ 随标的价格增加、随执行价格减少；期限延长扩大了可选择的执行时点，故美式期权价值不减；无套利给出 $c\ge\max(S_0-Ke^{-rT},0)$ 和 $p\ge\max(Ke^{-rT}-S_0,0)$。若任一关系违反，则买入支付较高而价格较低的组合、卖出支付较低而价格较高的组合即可构造套利。
\end{solution}
\begin{problem}
\problemno{6.9} 假设 $c_1$、$c_2$ 和 $c_3$ 分别表示执行价格分别为 $X_1$、$X_2$ 和 $X_3$ 的欧式看涨期权的价格，其中 $X_1<X_2<X_3$，并且 $X_3-X_2=X_2-X_1$。三个期权的标的资产相同、到期日相同。证明：$2c_2\le c_1+c_3$。
\end{problem}
\begin{solution}
对任意到期价格 $S_T$，函数 $c(K)=(S_T-K)^+$ 是执行价格 $K$ 的凸函数。由于 $X_2=(X_1+X_3)/2$，有
\[
2(S_T-X_2)^+\le(S_T-X_1)^++(S_T-X_3)^+.
\]
对状态价格加权并求和，得到 $2c_2\le c_1+c_3$。若不成立，则买入右侧较便宜的一方并卖出左侧较昂贵的一方即可套利。
\end{solution}
\begin{problem}
\problemno{6.10} 蝶形看涨价差是由标的资产相同、到期日相同而执行价格不相同的欧式看涨期权形成的期权组合：(1) 买入一份执行价格为 $K-\delta$ 的看涨期权和一份执行价格为 $K+\delta$ 的看涨期权；(2) 卖出 $2$ 份执行价格为 $K$ 的看涨期权。请画出蝶形价差的损益图，并证明当 $K$ 给定而 $\delta$ 不同时，$\delta$ 越大，蝶形价差的价格越高。
\end{problem}
\begin{solution}
到期支付为
\[
B(S_T)=(S_T-K+\delta)^+-2(S_T-K)^++(S_T-K-\delta)^+.
\]
因此 $S_T\le K-\delta$ 或 $S_T\ge K+\delta$ 时支付为零；在 $K-\delta<S_T\le K$ 时支付为 $S_T-K+\delta$，在 $K<S_T<K+\delta$ 时支付为 $K+\delta-S_T$，图形是顶点为 $(K,\delta)$ 的三角形。由凸性不等式，$c(K-\delta)+c(K+\delta)-2c(K)\ge0$；增大 $\delta$ 扩大非负支付区域，且逐状态支付不减，故无套利价格随 $\delta$ 增大而上升。
\end{solution}
\begin{problem}
\problemno{6.11} 考虑一个二叉树模型，股票 A 的期初价格为 \num{10} 元，上涨和下跌等概率，且 $u=20\%$，$d=-20\%$，以 A 为标的欧式看涨期权，无风险利率为 $10\%$。
\begin{enumerate}[label=(\alph*)]
    \item 画出股票 A 和执行价格为 \num{9} 元的欧式看涨期权的二叉树图形，计算期权在期初的价格；
    \item 假如得到新的利好消息：股票上涨的概率大于 $0.5$，股票期初价格仍为 \num{10} 元。请问你愿意购买看涨期权还是看跌期权？为什么？
\end{enumerate}
\end{problem}
\begin{solution}
\begin{enumerate}[label=(\alph*)]
\item 一期股票价格为 $12$ 或 $8$，执行价为 $9$ 的看涨期权支付为 $3$ 或 $0$。风险中性概率满足 $1.2p^*+0.8(1-p^*)=1.1$，故 $p^*=3/4$，期权初始价格为 $c_0=(3/4)3/1.1=45/22\approx2.045$。
\item 利好消息提高上涨的主观概率时，看涨期权的主观期望收益增加；若投资者相信该消息，应选择看涨期权。风险中性概率用于无套利定价，不等于真实概率。
\end{enumerate}
\end{solution}
\begin{problem}
\problemno{6.12} 在第 6.11 题的基础上考虑三时点两期的二叉树模型，上涨和下跌率保持不变，无风险利率也不变。
\begin{enumerate}[label=(\alph*)]
    \item 画出标的股票和欧式看涨期权的二叉树图形，并计算各节点的期权价格；
    \item 若期权为美式看涨期权，画出看涨期权的二叉树图形，计算各节点的期权价格。
\end{enumerate}
\end{problem}
\begin{solution}
取风险中性概率 $p^*=1/2$。四条路径的期末支付依次为 $7,1,0,0$，因为对应平均价格为 $12,6,3,1.5$。因此期初无套利价格为
\[
c_0=\frac{(7+1)/4}{1.25^2}=1.28。
\]
\end{solution}
\begin{problem}
\problemno{6.13} 同样考虑三时点(0、1、2 时点)两期二叉树模型，标的资产期初价格为 $S_0=4$ 元，$u=2$，$d=1/2$，无风险利率为 $25\%$。以该股票为标的资产的亚式看涨期权，期权最终价值为 $c_2=\max\bigl[(S_1+S_2)/2-5,0\bigr]$，即执行价格为 \num{5} 元，以股票在时点 1 和 2 处的价格 $S_1$、$S_2$ 的算术平均值来结算。请计算期权的期初无套利价格。
\end{problem}
\begin{solution}
取 $p^*=1/2$，四条路径的 $(S_1,S_2)$ 为 $(8,16),(8,4),(2,4),(2,1)$，平均价格为 $12,6,3,1.5$，期末支付为 $7,1,0,0$。因每条路径概率为 $1/4$，
\[
c_0=\frac{(7+1)/4}{(1.25)^2}=1.28。
\]
\end{solution}
\begin{problem}
\problemno{6.14} 假定期末有 $4$ 个不同的状态发生，假设股票 A 的支付向量为 $(1,2,3,4)^{\mathrm{T}}$，请你以股票 A 为基础构造欧式看涨期权使市场完全化。用记号表示期权和股票的价格，用这些符号表示出状态价格。
\end{problem}
\begin{solution}
选取执行价格 $1,2,3$ 的三个欧式看涨期权，连同股票 A 形成四个证券。支付向量为
\[
A=(1,2,3,4),\ C_1=(0,1,2,3),\ C_2=(0,0,1,2),\ C_3=(0,0,0,1)。
\]
其支付矩阵为上三角矩阵且行列式为 $1$，所以市场完全。若股票和期权价格分别为 $S,c_1,c_2,c_3$，状态价格满足
\[
q_4=c_3,\quad q_3=c_2-2c_3,\quad q_2=c_1-2c_2+c_3,\quad q_1=S-2c_1+c_2。
\]
\end{solution}
